% This must be in the first 5 lines to tell arXiv to use pdfLaTeX, which is strongly recommended.
\pdfoutput=1
% In particular, the hyperref package requires pdfLaTeX in order to break URLs across lines.
\documentclass[11pt]{article}
\usepackage{amsmath}
\usepackage{tabularx} % Add this in the preamble
\usepackage{amssymb}
% Remove the "review" option to generate the final version.
\usepackage{ACL2023}

% Standard package includes
\usepackage{times}
\usepackage{latexsym}
\usepackage{algorithm}
\usepackage{algorithmic}
\usepackage{graphicx}
\usepackage{subcaption}
\usepackage{placeins}
\usepackage{float}

% For proper rendering and hyphenation of words containing Latin characters (including in bib files)
\usepackage[T1]{fontenc}
% For Vietnamese characters
% \usepackage[T5]{fontenc}
% See https://www.latex-project.org/help/documentation/encguide.pdf for other character sets

% This assumes your files are encoded as UTF8
\usepackage[utf8]{inputenc}

% This is not strictly necessary, and may be commented out.
% However, it will improve the layout of the manuscript,
% and will typically save some space.
\usepackage{microtype}

% This is also not strictly necessary, and may be commented out.
% However, it will improve the aesthetics of text in
% the typewriter font.
\usepackage{inconsolata}
\usepackage{graphicx} 
\usepackage{float}
\usepackage{colortbl}
\usepackage{xcolor}
\usepackage{booktabs}
\usepackage{multirow}

% If the title and author information does not fit in the area allocated, uncomment the following
%
%\setlength\titlebox{<dim>}
%
% and set <dim> to something 5cm or larger.

\begin{document}

\onecolumn

\begin{center}
{\Large\bfseries Pixel-level Certified Explanations via Randomized Smoothing
\\ Application to Medical Imaging Classifiers Project Report \par}
\author{Daniya Niazi \\
\texttt{7062343,dani00003@stud.uni-saarland.de} \\
Universit{\"a}t des Saarlandes}
\end{center}

\vspace{1em}
\section{Introduction}
This report documents the implementation of pixel-level certified explanations for medical imaging classifiers, based on \cite{anani2025pixel} and adapted to four datasets (ChestX-ray14, Brain MRI, APTOS, ISIC). The project focuses on producing attribution maps with provable per-pixel robustness guarantees using randomized smoothing and top-$K$ sparsification. In addition to standard medical datasets, a grid-based evaluation (Certified GridPG) is implemented to test whether certified important pixels align with a known target region.

The goals are: (i) build a unified attribution pipeline across methods, (ii) certify attributions at the pixel level under $\ell_2$-bounded input perturbations, (iii) evaluate robustness and localization, and (iv) document deviations from the original paper introduced by practical implementation choices.

\section{Background}
\subsection{Attribution Methods and Their Fragility}
Post-hoc attribution methods map an image to a heatmap of pixel importance for a predicted class. This work employs gradient-based (Integrated Gradients), activation-based (Grad-CAM), and perturbation-based (RISE, Occlusion) methods, plus LRP. While these methods are widely used, their outputs can be unstable: small, visually imperceptible perturbations can significantly change the heatmap without changing the prediction. This fragility motivates the need for certified explanations that provide robustness guarantees, not just qualitative visualization.
\subsection{XAI in Medical Imaging: Necessity and
Limitations}
Medical decisions depend on localized visual evidence (lesions, opacities, tumors), so attribution maps are often used for sanity checks and clinical interpretability. However, medical images are high-stakes and sensitive to noise, artifacts, and domain shift. Unstable explanations can mislead clinicians or hide spurious correlations. Certified attributions aim to reduce this risk by stating, for each pixel, whether its importance label is robust within a provable radius.


\section{Pixel-Level Certification for
Attributions on Medical Dataset}

\subsection{Overview}
The paper \cite{anani2025pixel} pipeline is implemented in two settings:
(i) standard medical datasets (ISIC, ChestX-ray14, Brain MRI, APTOS) using single-head classifiers, and
(ii) a grid-based setting for Certified GridPG, where each grid cell contains a different class image and the target cell is known. In both cases, attribution maps are computed, sparsified to top-$K$ pixels, and randomized smoothing is applied to certify the per-pixel labels.

The implementation follows a unified workflow: compute attributions, sparsify to top-$K$, apply randomized smoothing, and output a ternary certified map with labels 1 (important), 0 (unimportant), or abstain.

\subsection{Paper Implementation and Algorithm}

This section presents the complete certification workflow with formal algorithms for each stage: attribution computation, top-$K$ sparsification, randomized smoothing, and final certification.

\subsubsection{Overview}

Given image $x$, the pipeline: (1) compute attribution heatmap, (2) keep top $K\%$ pixels, (3) add noise and aggregate, (4) certify pixels as important/unimportant/abstain. Any perturbation $\|\delta\|_2 \le R = \sigma \Phi^{-1}(\tau)$ preserves certified labels.

\subsubsection{Algorithm 1: Top-$K$ Sparsification}

\begin{algorithm}[H]
\caption{Top-$K$ Sparsification}
\label{alg:topk_sparsification}
\begin{algorithmic}[1]
\REQUIRE Heatmap $h(x)$, sparsity $K\%$
\ENSURE Binary map $h_K(x) \in \{0,1\}$
\STATE Rank all pixels by attribution value
\STATE Mark top $K\%$ as 1, rest as 0
\RETURN $h_K(x)$
\end{algorithmic}
\end{algorithm}

\subsubsection{Algorithm 2: Randomized Smoothing}

\begin{algorithm}[H]
\caption{Randomized Smoothing}
\label{alg:randomized_smoothing}
\begin{algorithmic}[1]
\REQUIRE Image $x$, method $h$, $K$, $\sigma$, samples $n=100$
\ENSURE Probabilities $p_1(i), p_0(i)$ per pixel
\FOR{$j = 1$ to $n$}
    \STATE $x^{(j)} \leftarrow x + \epsilon$, $\epsilon \sim \mathcal{N}(0, \sigma^2)$
    \STATE $h_K(x^{(j)}) \leftarrow$ TopK$(h(x^{(j)}), K)$
    \STATE Count important/unimportant labels per pixel
\ENDFOR
\STATE $p_1(i) = \text{count}_1[i] / n$, $p_0(i) = \text{count}_0[i] / n$
\RETURN $p_1, p_0$
\end{algorithmic}
\end{algorithm}

\subsubsection{Algorithm 3: Certification}

\begin{algorithm}[H]
\caption{Certification}
\label{alg:certification}
\begin{algorithmic}[1]
\REQUIRE Probabilities $p_1(i), p_0(i)$, threshold $\tau=0.75$
\ENSURE Certified map $c(x)$
\FOR{each pixel $i$}
    \IF{$p_1(i) > \tau$} $c(x)_i \leftarrow 1$
    \ELSIF{$p_0(i) > \tau$} $c(x)_i \leftarrow 0$
    \ELSE $c(x)_i \leftarrow$ abstain
    \ENDIF
\ENDFOR
\RETURN $c(x)$
\end{algorithmic}
\end{algorithm}


All experiments use: $\sigma = 0.15$, $n = 100$, $\tau = 0.75$, $K \in \{5\%, 25\%, 50\%\}$, certified radius $R \approx 0.101$.


\section{Implementation}

\subsection{Dataset Details and Characteristics}

This section describes the medical imaging datasets used in this seminar project for evaluating pixel-level certified explanations. The datasets span multiple imaging modalities and clinical tasks, enabling an assessment of certified attribution methods across diverse visual characteristics, class structures, and safety-critical scenarios. For each dataset, the imaging modality, classification task, and clinical relevance are outlined, with representative sample images provided.

\subsubsection{ChestX-ray14}

The ChestX-ray14 dataset consists of frontal chest radiographs acquired using X-ray transmission imaging. In this project, the task is formulated as a binary classification problem distinguishing pneumonia from normal cases. Chest X-rays are among the most frequently used diagnostic imaging tools worldwide, making them a high-impact and safety-critical application domain. Representative samples from the dataset are shown in Figure~\ref{fig:chestxray}.

\begin{figure}[H]
    \centering
    \includegraphics[width=1\linewidth]{data sample/chestsample.png}
    \caption{Representative samples from the ChestX-ray14 dataset showing normal and pneumonia cases.}
    \label{fig:chestxray}
\end{figure}

\subsubsection{Brain MRI Tumor Dataset}

The Brain MRI Tumor dataset consists of T1-weighted magnetic resonance images, released under a Creative Commons license. The dataset is pre-cleaned, resized, and augmented, and supports a four-class classification task: glioma, meningioma, pituitary tumor, and normal brain tissue. The class distribution in the training set is approximately balanced. Representative samples are shown in Figure~\ref{fig:brainmri}.

\begin{figure}[H]
    \centering
    \includegraphics[width=1\linewidth]{data sample/brain sample.png}
    \caption{Sample images from the Brain MRI Tumor dataset illustrating different tumor types and normal anatomy.}
    \label{fig:brainmri}
\end{figure}

\subsubsection{APTOS 2019 (Diabetic Retinopathy)}

The APTOS 2019 dataset consists of color fundus photographs captured using retinal imaging devices. The task is a five-class ordinal classification corresponding to diabetic retinopathy (DR) severity: Grade 0 (no DR), Grade 1 (mild), Grade 2 (moderate), Grade 3 (severe), and Grade 4 (proliferative).

The training data reflects real-world clinical prevalence, with approximately 49\% normal cases, followed by moderate (20\%), mild (13\%), severe (11\%), and proliferative (7\%) DR. Diabetic retinopathy is a leading cause of blindness, and early detection depends on identifying subtle vascular and color-based features. This makes fundus images particularly suitable for analyzing how certified explanations behave with respect to color and fine-grained texture cues. Representative samples are shown in Figure~\ref{fig:aptos}.

\begin{figure}[H]
    \centering
    \includegraphics[width=1\linewidth]{data sample/fundus.png}
    \caption{Representative fundus images from the APTOS 2019 dataset across different diabetic retinopathy severity levels.}
    \label{fig:aptos}
\end{figure}

\subsubsection{ISIC Skin Lesion Dataset}

The ISIC Skin Lesion dataset contains dermoscopic images captured using magnified color photography. The dataset supports a seven-class classification task: melanoma, melanocytic nevus, basal cell carcinoma, actinic keratosis, benign keratosis, dermatofibroma, and vascular lesions.

The dataset exhibits significant class imbalance, with melanocytic nevi comprising approximately 60\% of samples, while melanoma accounts for roughly 10\%. Dermoscopic images present substantial visual complexity due to color variation, texture patterns, and common artifacts such as hair or rulers. To ensure computational feasibility and balanced evaluation, the dataset was subsampled to a maximum of 500 images per class. Representative examples are shown in Figure~\ref{fig:isic}.

\begin{figure}[H]
    \centering
    \includegraphics[width=1\linewidth]{data sample/isic sample.png}
    \caption{Example dermoscopic images from the ISIC dataset illustrating diverse lesion types and visual variability.}
    \label{fig:isic}
\end{figure}

\subsubsection{Dataset Configuration}

All datasets were split into training, validation, and test sets using an 80/10/10 ratio. Images from all datasets were resized to a uniform resolution of $224 \times 224$ pixels to ensure consistency across models and attribution methods.

\begin{table}[h!]
\centering
\caption{Summary of datasets used in the certification experiments.}
\begin{tabularx}{\linewidth}{|X|X|X|X|X|}
\hline
\textbf{Dataset} & \textbf{Image Modality} & \textbf{Train-Test-Val Split} & \textbf{Image Size} \\
\hline
ChestX-ray14 & X-ray (grayscale) & 80-10-10\% & $224 \times 224$ \\
\hline
Brain MRI Tumor & MRI & 80-10-10\% & $224 \times 224$ \\
\hline
APTOS 2019 & Fundus photography& 80-10-10\% & $224 \times 224$ \\
\hline
ISIC Skin Lesion & Dermoscopic images & 80-10-10\% & $224 \times 224$ \\
\hline
\end{tabularx}
\label{tab:dataset_summary}
\end{table}


\subsection{Model Architecture Selection and Training Results}

This section presents a consolidated comparison of convolutional neural network (CNN) architectures evaluated across multiple medical imaging datasets. Instead of reporting dataset-specific results in separate paragraphs, validation accuracy is summarized in a single table to facilitate direct comparison across models and modalities. This analysis informs the selection of representative architectures for subsequent certified attribution experiments.

\subsubsection{Model Architectures}

Six CNN architectures were evaluated: ResNet-18, ResNet-50, DenseNet-121, EfficientNet-B0, EfficientNet-B1, and MobileNetV2. All models were initialized with ImageNet-1k pre-trained weights and fine-tuned on the medical datasets using the Adam optimizer with a learning rate of $1 \times 10^{-4}$, batch size 32, and early stopping with a patience of five epochs. Standard data augmentation was applied, including rotations ($\pm15^\circ$), horizontal flipping, and color jitter for RGB datasets only.

\subsubsection{Validation Accuracy Across Datasets}

Table~\ref{tab:val_acc_models} reports the validation accuracy of each model across the four medical imaging datasets. For each dataset, the highest validation accuracy is highlighted. While EfficientNet-B1 consistently achieves the strongest predictive performance, DenseNet-121 and ResNet-based architectures remain competitive. Simpler models such as ResNet-18 and MobileNetV2 are later favored for detailed certification analysis due to their computational efficiency and more stable attribution behavior.
\begin{table}[H]
\centering
\caption{Validation accuracy (\%) of CNN architectures across medical imaging datasets. Best-performing model per dataset is highlighted.}
\small
\setlength{\tabcolsep}{4pt} % default is ~6pt
\renewcommand{\arraystretch}{1.15}
\begin{tabularx}{\linewidth}{|X|c|c|c|c|c|c|}
\hline
\textbf{Dataset} 
& \textbf{ResNet-18} 
& \textbf{ResNet-50} 
& \textbf{DenseNet-121} 
& \textbf{EfficientNet-B0} 
& \textbf{EfficientNet-B1} 
& \textbf{MobileNetV2} \\
\hline
ChestX-ray14 
& 88.4 & 89.1 & 89.8 & 89.4 & \cellcolor{gray!25}\textbf{91} & 87.6 \\
\hline
Brain MRI Tumor 
& 90.9 & 91.7 & 92.3 & 92.0 & \cellcolor{gray!25}\textbf{92.7} & 90.1 \\
\hline
APTOS 2019 (DR) 
& 86.8 & 87.6 & 88.1 & 87.9 & \cellcolor{gray!25}\textbf{88.7} & 86.2 \\
\hline
ISIC Skin Lesion 
& 84.4 & 85.5 & 86.2 & 85.9 & \cellcolor{gray!25}\textbf{86.9} & 83.8 \\
\hline
\end{tabularx}
\label{tab:val_acc_models}
\end{table}


\begin{figure}[H]
\centering
\includegraphics[width=1\linewidth]{confusion matrix/brain_mri_confusion_matrix.png}
\caption{Confusion Matrix - Brain MRI}
\end{figure}

\begin{figure}[H]
\centering
\includegraphics[width=1\linewidth]{confusion matrix/chestxray_confusion_matrix.png}
\caption{Confusion Matrix - Chest X-ray}
\end{figure}

\begin{figure}[H]
\centering
\includegraphics[width=1\linewidth]{confusion matrix/confusion_matrix.png}
\caption{Confusion Matrix - ISIC}
\end{figure}

\begin{figure}[H]
\centering
\includegraphics[width=1\linewidth]{confusion matrix/fundus_confusion_matrix.png}
\caption{Confusion Matrix - APTOS Fundus}
\end{figure}



\subsubsection{Model Selection for Certification Experiments}

Although EfficientNet-B1 achieves the highest validation accuracy across all datasets, ResNet-18 was selected as the primary model for extensive certification experiments (Experiment 2). This choice reflects a trade-off between predictive performance, computational efficiency, and attribution interpretability, as illustrated on representative samples (e.g., 10 ISIC images shown in Figure~\ref{fig:isic}). In addition, ResNet-18, MobileNetV2, DenseNet-121, and EfficientNet-B1 were evaluated on a limited subset of samples ($N=5$ per dataset) to analyze architectural differences in certified attribution behavior.


\newpage
\subsection{Attribution Methods: Settings and Implementation}

This section describes the attribution methods used in this work. For each method, the intuition is briefly explained, a simplified pseudo-algorithm is presented, and the concrete hyperparameter settings used in the implementation are described. All methods produce pixel-level heatmaps compatible with the certification framework.

% ---------------------------------------------------------
\subsubsection{Integrated Gradients}
% ---------------------------------------------------------

\paragraph{Method intuition.}
Integrated Gradients (IG) estimates pixel importance by averaging gradients along a straight-line path from a baseline image to the input image, reducing sensitivity to local gradient noise.

\begin{algorithm}[H]
\caption{Integrated Gradients (IG)}
\label{alg:ig}
\begin{algorithmic}[1]
\REQUIRE Model $f_\theta$, input $x$, baseline $x'$, target class $c$, steps $N$
\ENSURE Attribution map $h(x)$
\FOR{$\alpha = 0$ to $1$ in $N$ steps}
    \STATE $x_\alpha \leftarrow x' + \alpha (x - x')$
    \STATE Compute gradient $\nabla_x f_\theta(x_\alpha)_c$
\ENDFOR
\STATE Average gradients and compute $(x - x') \odot \nabla_x$
\STATE Sum across channels and normalize
\RETURN $h(x)$
\end{algorithmic}
\end{algorithm}

\paragraph{Implementation details.}
Integrated Gradients uses a zero-valued (black) image as the baseline and approximates the path integral using 50 uniformly spaced interpolation steps. Gradients are accumulated along the interpolation path, multiplied by the input difference, and summed across channels to obtain a pixel-level heatmap. The final attribution map is normalized to the range $[0,1]$ for visualization and certification compatibility.

\paragraph{Remarks.}
Integrated Gradients produces smooth and spatially coherent attribution maps due to gradient averaging along the interpolation path. This implicit smoothing improves robustness to small input perturbations, which explains its consistently strong certification performance across datasets. However, the method requires multiple forward and backward passes per image, making it more computationally expensive than single-pass gradient methods.

% ---------------------------------------------------------
\subsubsection{Grad-CAM}
% ---------------------------------------------------------

\paragraph{Method intuition.}
Grad-CAM highlights spatial regions by weighting convolutional feature maps using gradients of the target class, producing coarse but stable explanations.

\begin{algorithm}[H]
\caption{Grad-CAM}
\label{alg:gradcam}
\begin{algorithmic}[1]
\REQUIRE Model $f_\theta$, target conv layer $L$, input $x$, class $c$
\ENSURE Heatmap $h(x)$
\STATE Store activations $A^k$ from layer $L$
\STATE Compute gradients $\partial f_\theta(x)_c / \partial A^k$
\STATE Compute weights $\alpha_k = \text{mean}(\nabla A^k)$
\STATE $h(x) = \sum_k \alpha_k A^k$
\STATE Apply ReLU and upsample
\RETURN $h(x)$
\end{algorithmic}
\end{algorithm}

\paragraph{Implementation details.}
In our implementation, Grad-CAM targets the final convolutional block (\texttt{layer4}) of the network. The resulting low-resolution activation map is upsampled to $224 \times 224$ using bilinear interpolation and passed through a ReLU operation to remove negative contributions. This produces a coarse but spatially stable attribution map aligned with the input image.

\paragraph{Remarks.}
Grad-CAM provides stable but coarse explanations, as relevance is derived from low-resolution convolutional feature maps. The inherent spatial downsampling acts as an implicit regularizer, improving robustness under certification. However, this comes at the cost of localization precision, particularly for small or fine-grained structures in medical images.

% ---------------------------------------------------------
\subsubsection{Layer-wise Relevance Propagation (LRP)}
% ---------------------------------------------------------

\paragraph{Method intuition.}
LRP redistributes the prediction score backward through the network while conserving total relevance, yielding sharp and localized explanations.

\begin{algorithm}[H]
\caption{Layer-wise Relevance Propagation (LRP)}
\label{alg:lrp}
\begin{algorithmic}[1]
\REQUIRE Model $f_\theta$, input $x$, target class $c$
\ENSURE Relevance heatmap $h(x)$
\STATE Forward pass to obtain prediction
\STATE Initialize relevance at output neuron
\FOR{each layer from output to input}
    \STATE Redistribute relevance using $\epsilon$-rule
\ENDFOR
\STATE Sum relevance across channels
\RETURN $h(x)$
\end{algorithmic}
\end{algorithm}

\paragraph{Implementation details.}
LRP is implemented using the Zennit framework with the $\epsilon$-rule ($\epsilon = 10^{-6}$) to ensure numerical stability during relevance propagation. Channel-wise relevance maps are summed and converted to absolute values to obtain pixel-level attributions. Importantly, normalization is applied only after certification to preserve the relative ranking of relevance scores.

\paragraph{Remarks.}
LRP generates sharp and highly localized attribution maps, often aligning well with visually salient regions. However, relevance propagation is sensitive to activation discontinuities (e.g., ReLU), leading to high variability under small perturbations. Without explicit robustness constraints, this sensitivity limits certification performance despite strong qualitative results on selected samples.

% ---------------------------------------------------------
\subsubsection{RISE}
% ---------------------------------------------------------

\paragraph{Method intuition.}
RISE estimates pixel importance by randomly masking parts of the image and measuring how the model’s prediction changes.

\begin{algorithm}[H]
\caption{RISE}
\label{alg:rise}
\begin{algorithmic}[1]
\REQUIRE Model $f_\theta$, input $x$, class $c$, masks $N$
\ENSURE Heatmap $h(x)$
\FOR{$i = 1$ to $N$}
    \STATE Sample random mask $M_i$
    \STATE Compute logit $f_\theta(x \odot M_i)_c$
    \STATE Accumulate $f_\theta(x \odot M_i)_c \cdot M_i$
\ENDFOR
\STATE Normalize attribution map
\RETURN $h(x)$
\end{algorithmic}
\end{algorithm}

\paragraph{Implementation details.}
RISE generates 500 random binary masks at a coarse resolution of $7 \times 7$, which are upsampled to the input image size using bilinear interpolation. Each masked image is forwarded through the model, and the target class logit is used as the weighting score for attribution accumulation. The final attribution map is normalized after aggregating contributions from all masks.

\paragraph{Remarks.}
Due to its stochastic nature, RISE exhibits high variance. Under certification, this variance is amplified, resulting in low robustness (20--28\%). The method is also computationally expensive due to the large number of forward passes.

% ---------------------------------------------------------
\subsubsection{Occlusion}
% ---------------------------------------------------------

\paragraph{Method intuition.}
Occlusion measures importance by systematically hiding image patches and observing the drop in prediction confidence.

\begin{algorithm}[H]
\caption{Occlusion}
\label{alg:occlusion}
\begin{algorithmic}[1]
\REQUIRE Model $f_\theta$, input $x$, class $c$
\ENSURE Heatmap $h(x)$
\FOR{each patch location}
    \STATE Replace patch with baseline value
    \STATE Measure logit drop $\Delta$
    \STATE Assign $\Delta$ to patch
\ENDFOR
\STATE Upsample attribution map
\RETURN $h(x)$
\end{algorithmic}
\end{algorithm}

\paragraph{Implementation details.}
Occlusion uses square patches of size $16 \times 16$ with a stride of 8 pixels, resulting in overlapping occlusions across the image. Occluded regions are replaced with the mean pixel value of the image to avoid introducing artificial artifacts. The resulting low-resolution attribution map is then upsampled to full image resolution using bilinear interpolation.

\paragraph{Remarks.}
Occlusion-based methods are highly sensitive to patch alignment. Small input shifts change which regions are occluded, leading to unstable attribution maps. Combined with high computational cost, this explains the low certification rates (12--22\%) observed.



% =========================================================
\subsection{Attribution Visualization Across Datasets}
% =========================================================

This section presents qualitative attribution visualizations across datasets. For each dataset, two representative samples are selected from the ResNet-18 model. Attribution maps are generated using Integrated Gradients, Grad-CAM, RISE, Occlusion, and LRP.

\subsubsection{Brain MRI Tumor}
\begin{figure}[H]
\centering
\includegraphics[width=1\linewidth]{attribution sample/brainatt1resnet18.png}
\caption{Attribution visualizations for Brain MRI Sample 1 using ResNet-18. From left to right: input image, Integrated Gradients, Grad-CAM, RISE, Occlusion, and LRP.}
\end{figure}

\begin{figure}[H]
\centering
\includegraphics[width=1\linewidth]{attribution sample/brainattr2resnet18.png}
\caption{Attribution visualizations for Brain MRI Sample 2 using ResNet-18. From left to right: input image, Integrated Gradients, Grad-CAM, RISE, Occlusion, and LRP.}
\end{figure}



\subsubsection{Chest X-ray}

\begin{figure}[H]
\centering
\includegraphics[width=1\linewidth]{attribution sample/chestattrsample1.png}
\caption{Attribution visualizations for Chest X-ray Sample 1 using ResNet-18.}
\end{figure}

\begin{figure}[H]
\centering
\includegraphics[width=1\linewidth]{attribution sample/chestattrsample2.png}
\caption{Attribution visualizations for Chest X-ray Sample 2 using ResNet-18.}
\end{figure}


\subsubsection{APTOS 2019 (DR)}

\begin{figure}[H]
\centering
\includegraphics[width=1\linewidth]{attribution sample/aptosatt1sample1.png}
\caption{Attribution visualizations for APTOS 2019 (DR) Sample 1 using ResNet-18.}
\end{figure}

\begin{figure}[H]
\centering
\includegraphics[width=1\linewidth]{attribution sample/aptosatt1sample2.png}
\caption{Attribution visualizations for APTOS 2019 (DR) Sample 2 using ResNet-18.}
\end{figure}


\subsubsection{ISIC Skin Lesion }

\begin{figure}[H]
\centering
\includegraphics[width=1\linewidth]{attribution sample/skinattr1.png}
\caption{Attribution visualizations for ISIC Skin Lesion Sample 1 using ResNet-18.}
\end{figure}

\begin{figure}[H]
\centering
\includegraphics[width=1\linewidth]{attribution sample/skinsample2.png}
\caption{Attribution visualizations for ISIC Skin Lesion Sample 2 using ResNet-18.}
\end{figure}


\subsection{Certification Visualization Across Datasets with ResNet-18}

This section presents certified attribution visualizations across all datasets. For each dataset, two representative validation samples are selected using the ResNet-18 model. Attribution maps are generated using Integrated Gradients, Grad-CAM, RISE, Occlusion, and LRP, and subsequently evaluated under a unified robustness certification framework.

\subsubsection{Standard Certification Framework and Hyperparameter Settings}

All certification experiments follow a consistent randomized smoothing protocol across datasets and attribution methods to ensure fair and comparable evaluation. Given an input image $x$ and its attribution heatmap $h(x)$, robustness is assessed by injecting Gaussian noise into the input and analyzing the stability of the resulting attribution maps.

Specifically, for each image, we generate $N = 100$ noisy samples of the form
\[
x^{(i)} = x + \epsilon^{(i)}, \quad \epsilon^{(i)} \sim \mathcal{N}(0, \sigma^2 I),
\]
where $\sigma = 0.15$ is fixed for all experiments. For each noisy sample, the attribution method is recomputed, producing a set of noisy attribution maps that are subsequently aggregated into a smoothed sparsified (SS) attribution using the Top-K sparsification procedure described in Algorithm~\ref{alg:topk_sparsification}.

Certification is performed following the randomized smoothing approach (Algorithm~\ref{alg:randomized_smoothing}) by thresholding the smoothed attribution according to a confidence parameter $\tau = 0.75$. A pixel is considered \emph{certified} if its attribution label remains consistently above the threshold across noisy samples within the certified radius induced by $\sigma$, as formalized in Algorithm~\ref{alg:certification}. Pixels that do not meet this consistency requirement are marked as abstained. Certification is evaluated at three sparsity levels corresponding to the top $K=\{5\%, 25\%, 50\%\}$ most salient pixels.

The same certification parameters ($\sigma$, $N$, $\tau$, and $K$ values) are applied uniformly across all datasets and models. This standardized setup ensures that observed differences in certification behavior are attributable to dataset characteristics or attribution methods rather than hyperparameter tuning.


\subsubsection{Brain MRI Tumor}
\begin{figure}[H]
\centering
\includegraphics[width=1\linewidth]{certifcation/mri/resnet18_img5_paper_style.png}
\caption{Certification visualizations for Brain MRI Sample 1 using ResNet-18.}
\end{figure}

\begin{figure}[H]
\centering
\includegraphics[width=1\linewidth]{certifcation/mri/resnet18_img59_paper_style.png}
\caption{Certification visualizations for Brain MRI Sample 2 using ResNet-18.}
\end{figure}

\subsubsection{Chest X-ray}

\begin{figure}[H]
\centering
\includegraphics[width=1\linewidth]{certifcation/chest/resnet18_img117_paper_style.png}
\caption{Certification visualizations for Chest X-ray Sample 1 using ResNet-18.}
\end{figure}

\begin{figure}[H]
\centering
\includegraphics[width=1\linewidth]{certifcation/chest/resnet18_img31_paper_style.png}
\caption{Certification visualizations for Chest X-ray Sample 2 using ResNet-18.}
\end{figure}



\subsubsection{APTOS 2019 (DR)}

\begin{figure}[H]
\centering
\includegraphics[width=1\linewidth]{certifcation/fundus/resnet18_img146_paper_style.png}
\caption{Certification visualizations for APTOS 2019 (DR) Sample 1 using ResNet-18.}
\end{figure}

\begin{figure}[H]
\centering
\includegraphics[width=1\linewidth]{certifcation/fundus/resnet18_img65_paper_style.png}
\caption{Certification visualizations for APTOS 2019 (DR) Sample 2 using ResNet-18.}
\end{figure}



\subsubsection{ISIC Skin Lesion }

\begin{figure}[H]
\centering
\includegraphics[width=1\linewidth]{certifcation/isic/resnet18_img185_paper_style.png}
\caption{Certification visualizations for ISIC Skin Lesion Sample 1 using ResNet-18.}
\end{figure}

\begin{figure}[H]
\centering
\includegraphics[width=1\linewidth]{certifcation/isic/resnet18_img371_paper_style.png}
\caption{Certification visualizations for ISIC Skin Lesion Sample 2 using ResNet-18.}
\end{figure}


\subsection{Experiment 1 for Pixel-Level Certification: Cross-Architecture Comparison}

Experiment 1 evaluates all 4 architectures on a smaller subset (5 images per dataset, 20 total) using standard settings ($\sigma=1$, $N=50$, Integrated Gradients) to understand architecture effects on certification. 

\begin{table}[h!]
\centering
\caption{Experiment 1: Average certification accuracy (\%) across datasets and models for different top-$K$ certified pixel fractions.}
\label{tab:exp1_cert}
\resizebox{\textwidth}{!}{
\begin{tabular}{llccc ccc ccc ccc}
\toprule
& & \multicolumn{3}{c}{\textbf{Chest X-ray}} 
& \multicolumn{3}{c}{\textbf{Brain MRI}} 
& \multicolumn{3}{c}{\textbf{ISIC Skin}} 
& \multicolumn{3}{c}{\textbf{Fundus}} \\
\cmidrule(lr){3-5} \cmidrule(lr){6-8} \cmidrule(lr){9-11} \cmidrule(lr){12-14}
\textbf{Dataset} & \textbf{Model} 
& $K=5$ & $K=25$ & $K=50$
& $K=5$ & $K=25$ & $K=50$
& $K=5$ & $K=25$ & $K=50$
& $K=5$ & $K=25$ & $K=50$ \\
\midrule

\multirow{4}{*}{All} 
& ResNet-18     
& 92.90 & 49.71 & 36.51 
& 93.28 & 61.3 & 51.3 
& 89.84 & 45.36 & 31.97 
& 79.95 & 35.90 & 31.4 \\

& DenseNet-121  
& 91.24 & 35.81 & 25.00 
& 94.48 & 62.35 & 49.00 
& 92.20 & 41.57 & 30.26 
& 90.91 & 32.33 & 16.00 \\

& MobileNetV2   
& 89.87 & 40.97 & 24.37 
& 93.65 & 64.80 & 56.03 
& 84.97 & 44.85 & 33.38 
& 90.08 & 43.26 & 29.31 \\

& EfficientNet-B1
& 90.05 & 50.90 & 35.63 
& 92.13 & 55.62 & 42.68 
& 88.52 & 45.97 & 36.72 
& 89.73 & 42.37 & 26.47 \\

\bottomrule
\end{tabular}
}
\end{table}
As expected, certification accuracy decreases monotonically as the certified pixel fraction $K$ increases, reflecting the increasing difficulty of guaranteeing robustness for larger attribution regions.Among all, Brain MRI consistently achieves the highest certification rates, while fundus and ISIC exhibit lower robustness due to class difference, color and texture complexity. 


\begin{table}[h!]
\centering
\caption{Experiment 1: Average percentage of certified pixels across datasets (all models combined).}
\label{tab:exp1_dataset_avg}
\begin{tabular}{lcc}
\toprule
\textbf{Dataset} & \textbf{Mean Certified (\%)} & \textbf{Total Images} \\
\midrule
Brain MRI   & 67.96 & 300 \\
ISIC Skin  & 55.47 & 300 \\
Chest X-ray & 55.25 & 300 \\
Fundus     & 51.8 & 300 \\
\midrule
\textbf{Overall} & \textbf{57.31} & \textbf{1200} \\
\bottomrule
\end{tabular}
\end{table}

\begin{table}[h!]
\centering
\caption{Experiment 1: Average percentage of certified pixels per model and dataset.}
\label{tab:exp1_model_avg}
\resizebox{\textwidth}{!}{
\begin{tabular}{lcccc}
\toprule
\textbf{Dataset} & \textbf{ResNet-18} & \textbf{DenseNet-121} & \textbf{MobileNetV2} & \textbf{EfficientNet-B1} \\
\midrule
Brain MRI   & 68.28 & 68.61 & \textbf{71.49} & 63.47 \\
Chest X-ray & \textbf{59.70} & 51.8 & 51.74 & 58.86 \\
ISIC Skin  & 55.72 & 54.67 & 54.40 & \textbf{57.07} \\
Fundus     & 48.83 & 46.41 & \textbf{54.22} & 52.86 \\
\bottomrule
\end{tabular}
}
\end{table}



\subsection{Experiment 2: Large-Scale Certification with ResNet-18}

Experiment 2 extends the analysis of Experiment 1 by scaling the certification evaluation to a substantially larger number of samples while fixing the model architecture. Based on the stability, competitive performance, and computational efficiency observed in Experiment 1, \textbf{ResNet-18} was selected as the backbone model for all datasets.

While Experiment 1 evaluated certification behavior on a limited subset of images ($N=5$ per model–dataset pair), Experiment 2 increases the number of evaluated samples to up to $N=100$ images per dataset, resulting in a total of 6070 certified samples across all datasets. All certification hyperparameters, attribution methods, and noise settings are kept identical to Experiment 1 to ensure a controlled comparison.

\begin{table}[H]
\centering
\caption{Experiment 2: Average certification accuracy (\%) using ResNet-18 across datasets for different certified pixel fractions.}
\label{tab:exp2_kwise}
\begin{tabular}{lccc}
\toprule
\textbf{Dataset} & $K=5$ & $K=25$ & $K=50$ \\
\midrule
Brain MRI   & 93.49 & 61.31 & 49.57 \\
Chest X-ray & 92.26 & 49.37 & 36.80 \\
Fundus     & 76.69 & 36.19 & 29.83 \\
ISIC Skin  & 89.92 & 41.99 & 31.7 \\
\bottomrule
\end{tabular}
\end{table}

Consistent with Experiment 1, certification accuracy decreases monotonically as the certified pixel fraction $K$ increases. Brain MRI remains the most robust dataset, while Fundus and ISIC exhibit lower certification rates due to higher visual complexity and color variation.

\begin{table}[h!]
\centering
\caption{Experiment 2: Dataset-level average percentage of certified pixels using ResNet-18 across all attribution methods.}
\label{tab:exp2_dataset_avg}
\begin{tabular}{lc}
\toprule
\textbf{Dataset} & \textbf{Mean Certified (\%)} \\
\midrule
Brain MRI   & 68.12 \\
Chest X-ray & 59.48 \\
ISIC Skin  & 54.20 \\
Fundus     & 47.57 \\
\midrule
\textbf{Overall} & \textbf{57.31} \\
\bottomrule
\end{tabular}
\end{table}


\subsection{Evaluation: Robustness and Faithfulness}

This section presents a comprehensive evaluation of certified attributions along two complementary dimensions: robustness and faithfulness. Robustness quantifies the proportion of pixels whose attribution labels remain stable under input perturbations, while faithfulness measures whether certified important pixels actually contribute to the model's prediction.

\subsubsection{Robustness Via \%Certified}

Robustness is evaluated by measuring the percentage of pixels that receive a certified label (either important or unimportant) rather than abstaining. A higher certification rate indicates that the attribution method produces stable, predictable importance assignments that remain consistent under $\ell_2$-bounded perturbations within the certified radius $R = \sigma \Phi^{-1}(\tau) \approx 0.101$.
For each image, we compute:
$$
\text{Robustness} = \frac{\#\{\text{certified pixels (important or unimportant)}\}}{\text{total pixels}} \times 100\%
$$

Robustness decreases monotonically as the sparsity parameter $K$ increases, since maintaining consistent top-$K$ rankings becomes more difficult for larger $K$. Attribution methods with inherent smoothness (e.g., Integrated Gradients, Grad-CAM) achieve higher robustness than methods sensitive to local perturbations (e.g., RISE, Occlusion, LRP).

\subsubsection{Faithfulness Via Certified Pixel Deletion}

Faithfulness evaluates whether pixels certified as important actually influence the model's prediction. We measure this using a pixel deletion curve: certified important pixels are removed in order of decreasing attribution strength, and the model's confidence on the target class is tracked after each deletion step.
For a given sparsity level $K$, we:
\begin{enumerate}
\item Identify all pixels certified as important (label = 1).
\item Rank these pixels by their original attribution magnitude.
\item Sequentially replace the top $p\%$ of certified important pixels with a baseline value (mean pixel intensity).
\item Measure the remaining model confidence at each deletion fraction $p \in \{10\%, 20\%, ..., 100\%\}$.
\end{enumerate}

A steep drop in confidence indicates high faithfulness: removing certified important pixels significantly degrades the prediction, confirming that these pixels genuinely contribute to the model's decision. Conversely, a shallow curve suggests that certified pixels are not truly important.

\subsubsection{Experiment 1: Robustness}
\begin{figure}[H]
\centering
\includegraphics[width=1\linewidth]{experiment1/robustness/brain_mri/robustness_stacked.png}
\caption{Robustness certification results for Brain MRI across four ResNet-18 architectures for 5 images.}
\end{figure}

\begin{figure}[H]
\centering
\includegraphics[width=1\linewidth]{experiment1/robustness/chestxray/robustness_stacked.png}
\caption{Robustness certification results for Chest X-ray across four ResNet-18 architectures for 5 images.}
\end{figure}

\begin{figure}[H]
\centering
\includegraphics[width=1\linewidth]{experiment1/robustness/fundus/robustness_stacked.png}
\caption{Robustness certification results for APTOS 2019 across four ResNet-18 architectures for 5 images.}
\end{figure}

\begin{figure}[H]
\centering
\includegraphics[width=1\linewidth]{experiment1/robustness/isic/robustness_stacked.png}
\caption{Robustness certification results for ISIC Skin Lesion across four ResNet-18 architectures for 5 images.}
\end{figure}



\begin{figure}[H]
    \centering
    \includegraphics[width=1\linewidth]{experiment1/avg_robustness.png}
    \caption{Average Pixel Certified across dataset}
    \label{fig:chestxray}
\end{figure}

\paragraph{Robustness Insights from Cross-Architecture Comparison.}
The robustness evaluation in Experiment 1 reveals several key findings. First, \textbf{certification robustness decreases monotonically with increasing sparsity $K$} across all datasets, architectures, and attribution methods. This trend reflects the fundamental challenge of maintaining stable top-$K$ rankings under $\ell_2$-bounded perturbations: as $K$ increases, more pixels must be ranked consistently, making certification progressively harder.

Second, \textbf{gradient-based methods (IG, Grad-CAM) achieve substantially higher certification rates than perturbation-based methods (RISE, Occlusion, LRP)}. This superiority stems from their inherent smoothness: gradient-based attributions change continuously and predictably under input perturbations, whereas perturbation-based methods exhibit more erratic behavior due to their reliance on sampling or masking operations.

Third, \textbf{architecture choice has minimal impact on relative method rankings}. While absolute certification rates vary slightly across ResNet-18, DenseNet-121, MobileNetV2, and EfficientNet-B1, the robustness ordering among attribution methods remains consistent. This suggests that certification behavior is primarily determined by the attribution method's algorithmic properties rather than the underlying model architecture.

Finally, \textbf{dataset characteristics strongly influence certification difficulty}. Brain MRI, featuring high-contrast grayscale images with well-localized tumors, achieves the highest certification rates (68\%). In contrast, APTOS fundus images, characterized by complex color patterns and distributed vascular structures, exhibit the lowest certification rates (52\%). This demonstrates that medical imaging modality—specifically, color complexity and texture patterns—directly impacts attribution robustness.

\subsubsection{Experiment 1: Faithfullness}

\begin{figure}[H]
\centering
\includegraphics[width=1\linewidth]{experiment1/faithfullness/brain_mri/faithfulness_confidence_curves.png}
\caption{Faithfulness confidence curves for Brain MRI across four ResNet-18 architectures for 5 images at $K=\{5\%,25\%,50\%\}$.}
\end{figure}

\begin{figure}[H]
\centering
\includegraphics[width=1\linewidth]{experiment1/faithfullness/chestxray/faithfulness_confidence_curves.png}
\caption{Faithfulness confidence curves for Chest X-ray across four ResNet-18 architectures for 5 images at $K=\{5\%,25\%,50\%\}$.}
\end{figure}

\begin{figure}[H]
\centering
\includegraphics[width=1\linewidth]{experiment1/faithfullness/fundus/faithfulness_confidence_curves.png}
\caption{Faithfulness confidence curves for APTOS 2019 across four ResNet-18 architectures for 5 images at $K=\{5\%,25\%,50\%\}$.}
\end{figure}

\begin{figure}[H]
\centering
\includegraphics[width=1\linewidth]{experiment1/faithfullness/isic/faithfulness_confidence_curves.png}
\caption{Faithfulness confidence curves for ISIC Skin Lesion across four architectures for 5 validation images at $K=\{5\%,25\%,50\%\}$. Dermoscopic images exhibit varied faithfulness patterns across methods and architectures. Gradient-based methods demonstrate more consistent faithfulness, while perturbation-based methods show higher variance in confidence degradation.}
\label{fig:exp1_faithfulness_isic}
\end{figure}

\textbf{Average Faithfullness across dataset}


\begin{figure}[H]
    \centering
    \includegraphics[width=1\linewidth]{experiment1/faithfullness/overall_confidence_curves.png}
    \caption{Average confidence drop across all four datasets using ResNet-18 for Experiment 1. The curve aggregates faithfulness behavior across Brain MRI, Chest X-ray, APTOS, and ISIC. Steeper initial drop indicates that the most important certified pixels (first 10-30\% removed) are highly faithful, while the plateau suggests diminishing marginal contribution from lower-ranked certified pixels.}
    \label{fig:exp1_overall_faithfulness}
\end{figure}

\paragraph{Faithfulness Insights from Pixel Deletion Curves.}
The faithfulness evaluation complements robustness analysis by measuring whether certified important pixels genuinely contribute to model predictions. The pixel deletion curves reveal consistent patterns across datasets and architectures.

\textbf{Gradient-based methods demonstrate superior faithfulness.} IG and Grad-CAM consistently produce the steepest confidence drops when certified important pixels are removed, confirming that their top-ranked pixels are genuinely influential. Removing just 20-30\% of certified important pixels often reduces model confidence by 40-60\%, validating high faithfulness.

\textbf{Perturbation-based methods show shallower faithfulness curves.} RISE, Occlusion, and LRP exhibit more gradual confidence degradation, suggesting that their certified important pixels are less consistently aligned with the model's decision-making process. This faithfulness deficit partially explains why these methods also achieve lower robustness: unstable attributions are less likely to be both robust and faithful.

\textbf{Faithfulness degrades more slowly at higher $K$ values.} At $K=50\%$, confidence curves are shallower than at $K=5\%$, indicating that lower-ranked certified pixels contribute less to predictions. This aligns with the intuition that only the most salient pixels are critical for classification, while expanding to larger $K$ includes less influential regions.

\textbf{Architecture effects on faithfulness are minimal.} Similar to robustness, the four evaluated architectures (ResNet-18, DenseNet-121, MobileNetV2, EfficientNet-B1) exhibit nearly identical faithfulness patterns. This cross-architecture consistency reinforces that faithfulness is primarily an attribution method property rather than a model-specific characteristic.

\paragraph{Concerning Observations: Erratic Faithfulness Curves.}
Despite the overall trends described above, many individual faithfulness curves exhibit concerning patterns that warrant critical examination. Ideally, pixel deletion curves should decrease \emph{smoothly and monotonically}: as more certified important pixels are removed, model confidence should steadily decline. However, several observed curves show:

\begin{itemize}
\item \textbf{Flat or near-flat curves:} Removing up to 50-70\% of "certified important" pixels produces minimal confidence drop ($<10\%$). This suggests these pixels are not genuinely important to the model's decision, contradicting the attribution.
\item \textbf{Non-monotonic behavior:} Confidence \emph{increases} after deleting certain pixel fractions, then drops again. This counter-intuitive pattern indicates that some certified pixels may actually be \emph{suppressing} the prediction or that the deletion baseline introduces spurious signals.
\item \textbf{High inter-method variance:} On the same image, different attribution methods produce radically different curve shapes, suggesting attribution instability rather than robust signal.
\end{itemize}

These irregularities raise the possibility that randomized smoothing, rather than filtering unstable attributions, may be \emph{certifying noise}—pixel patterns that happen to be consistent across noisy samples by chance rather than reflecting genuine model reasoning. This concern is amplified by the grid localization failures (Section~\ref{tab:grid_isic_localization}), where certified pixels completely fail to align with known ground truth.

% EXPERIMENT 2


\subsubsection{Experiment 2: Robustness}
\begin{figure}[H]
\centering
\includegraphics[width=1\linewidth]{experiment2/robustness/mri/robustness_stacked.png}
\caption{Robustness certification results for Brain MRI using ResNet-18.}
\end{figure}

\begin{figure}[H]
\centering
\includegraphics[width=1\linewidth]{experiment2/robustness/chestxray/robustness_stacked.png}
\caption{Robustness certification results for Chest X-ray using ResNet-18.}
\end{figure}

\begin{figure}[H]
\centering
\includegraphics[width=1\linewidth]{experiment2/robustness/fundus/robustness_stacked.png}
\caption{Robustness certification results for APTOS 2019 using ResNet-18.}
\end{figure}

\begin{figure}[H]
\centering
\includegraphics[width=1\linewidth]{experiment2/robustness/isic/robustness_stacked.png}
\caption{Robustness certification results for ISIC Skin Lesion using ResNet-18.}
\end{figure}



\begin{figure}[H]
    \centering
    \includegraphics[width=1\linewidth]{experiment2/robustness/summary_mean_pct_certified.png}
    \caption{Average Pixel Certified across dataset}
    \label{fig:chestxray}
\end{figure}



\begin{figure}[H]
    \centering
    \includegraphics[width=1\linewidth]{experiment2/robustness/method_by_dataset.png}
    \caption{Average Pixel Certified across dataset - Attribution Method Di}
    \label{fig:chestxray}
\end{figure}

\subsubsection{Experiment 2: Faithfullness}

\begin{figure}[H]
\centering
\includegraphics[width=1\linewidth]{experiment2/faithfullness/mri/faithfulness_confidence_curves.png}
\caption{Faithfulness confidence curves for Brain MRI using ResNet-18 at $K=\{5\%,25\%,50\%\}$.}
\end{figure}

\begin{figure}[H]
\centering
\includegraphics[width=1\linewidth]{experiment2/faithfullness/chestxray/faithfulness_confidence_curves.png}
\caption{Faithfulness confidence curves for Chest X-ray using ResNet-18 at $K=\{5\%,25\%,50\%\}$.}
\end{figure}

\begin{figure}[H]
\centering
\includegraphics[width=1\linewidth]{experiment2/faithfullness/fundus/faithfulness_confidence_curves.png}
\caption{Faithfulness confidence curves for APTOS 2019 using ResNet-18 at $K=\{5\%,25\%,50\%\}$.}
\end{figure}

\begin{figure}[H]
\centering
\includegraphics[width=1\linewidth]{experiment2/faithfullness/isic/faithfulness_confidence_curves.png}
\caption{Faithfulness confidence curves for ISIC Skin Lesion using ResNet-18 at $K=\{5\%,25\%,50\%\}$.}
\end{figure}

\textbf{Average Faithfullness across dataset}


\begin{figure}[H]
    \centering
    \includegraphics[width=1\linewidth]{experiment2/faithfullness/overall_confidence_curves.png}
    \caption{Average Confidence drop across dataset with Resenet18 Model for experiment 2}
    \label{fig:chestxray}
\end{figure}

\paragraph{Scalability Validation and Large-Scale Insights.}
Experiment 2 extends the certification evaluation from 5 images per dataset to 100 images per dataset, scaling the total evaluation from 80 samples (Experiment 1: 4 architectures × 5 images × 4 datasets) to 400 samples (Experiment 2: 1 architecture × 100 images × 4 datasets). This large-scale evaluation validates the generalizability of findings from Experiment 1 and provides stronger statistical confidence in observed trends.

\textbf{Robustness patterns remain consistent at scale.} The dataset ranking observed in Experiment 1—Brain MRI (68\%) > Chest X-ray/ISIC (59\%, 54\%) > APTOS (48\%)—is preserved in Experiment 2, confirming that certification difficulty is inherently tied to dataset characteristics rather than sample-specific artifacts. The monotonic decrease in certification as $K$ increases continues to hold across all 400 samples, reinforcing the fundamental challenge of certifying larger top-$K$ sets.

\textbf{Method-wise performance differences become more pronounced.} With 100 images per dataset, the gap between gradient-based methods (IG, Grad-CAM) and perturbation-based methods (RISE, Occlusion, LRP) becomes more statistically significant. IG consistently achieves 10-20\% higher certification rates than RISE and Occlusion across all datasets, demonstrating that smoothness-induced robustness is not a small-sample phenomenon but a fundamental property.

\textbf{Faithfulness curves show reduced variance.} Aggregating over 100 images produces smoother confidence degradation curves compared to the 5-image evaluation. The steepness of initial drops (0-30\% deletion) remains consistent with Experiment 1, confirming that faithful attribution is a reproducible property across large sample sizes. IG and Grad-CAM maintain superior faithfulness, with confidence drops of 40-60\% after removing just 30\% of certified pixels.

\textbf{Color modality remains the primary challenge.} Both APTOS and ISIC (color datasets) exhibit systematically lower robustness and more variable faithfulness compared to Brain MRI and Chest X-ray (grayscale datasets). This pattern holds even when controlling for sample size, architecture, and hyperparameters, suggesting that color information introduces fundamental challenges for attribution stability. Texture complexity and multi-channel dependencies appear to amplify attribution noise under perturbations.

\paragraph{Persistent Faithfulness Concerns at Scale.}
Scaling to 100 images per dataset does not resolve the erratic faithfulness patterns observed in Experiment 1. While averaging across more samples produces visually smoother aggregate curves, inspection of individual image-level deletion curves reveals that the underlying issues persist:

\begin{itemize}
\item \textbf{Averaging masks individual failures:} Many images still exhibit flat or non-monotonic curves. The smooth aggregate curves in the figures result from averaging over diverse failure modes, not from consistent faithful behavior across images.
\item \textbf{RISE and Occlusion show minimal faithfulness impact:} Even after 80-90\% pixel deletion, confidence often remains above 50\%, indicating that these methods' certified pixels are largely irrelevant to predictions.
\item \textbf{LRP's sharp curves may reflect annotation artifacts:} LRP sometimes shows steep drops, but visual inspection suggests it may be identifying texture edges or color boundaries rather than diagnostically relevant regions.
\end{itemize}

Combined with the grid localization failures (where most methods achieve GridPG near or below the theoretical random expectation of 0.25), these observations suggest that the certification framework may be fundamentally limited when applied to medical images without stronger model training or attribution method redesign.


\paragraph{Effect of Increasing the Number of Samples.}
Increasing the number of evaluated samples in Experiment~2 did not lead to a qualitative change in the observed certification trends. The relative robustness ordering across datasets remained consistent with Experiment~1, with Brain MRI exhibiting the highest certification rates, followed by Chest X-ray and ISIC, and Fundus remaining the most challenging. Similarly, the monotonic decrease in certification as the certified pixel fraction $K$ increased was preserved. This indicates that the robustness patterns identified in the small-scale setting generalize to larger evaluations and are not artifacts of limited sample selection.



\section{Grid-Based Localization: Implementation Details}

This section presents the grid-based evaluation framework for testing whether certified attributions correctly localize to known target regions. We build $2\times2$ grids from ISIC images, placing a different class in each cell. The target cell is index 0 (top-left), and the target class is the class placed there. We generate 100 grid images per split and derive a binary target mask for the target cell. These grids are then passed through the same attribution and smoothing pipeline (Algorithms~1-3), and the certified maps are scored using the GridPG metric defined below.


\subsection{Grid Dataset Construction}

Grid images are synthetic $2 \times 2$ compositions of individual lesion images from different classes, as described in Algorithm~\ref{alg:grid_generation}. Each grid has a known target cell (top-left) whose class determines the grid's label. Images are sampled randomly from class pools, not based on model confidence. If a class pool is exhausted, it is refilled and reshuffled. Cell 0 (top-left) is always the target across all grids.

\begin{algorithm}[H]
\caption{Grid Dataset Generation}
\label{alg:grid_generation}
\begin{algorithmic}[1]
\REQUIRE Base dataset $D$ (ISIC), number of grids $N$, scale $m=2$, target cell index $t=0$
\ENSURE Grid dataset $\{G_i, y_i, M_i\}_{i=1}^N$ with images, labels, and ground-truth masks
\STATE \textbf{Organize base dataset by class}
\STATE Build class pools: $P_c \leftarrow \{i \mid D[i].\text{label} = c\}$ for each class $c$
\STATE Shuffle each pool: $P_c \leftarrow \text{shuffle}(P_c)$
\STATE \textbf{Generate each grid}
\FOR{$i = 1$ to $N$}
    \STATE \textbf{// Sample classes}
    \STATE $\mathcal{C} \leftarrow$ sample $m \times m$ distinct classes without replacement
    \STATE Randomly permute $\mathcal{C}$ and assign to cells: $\{c_0, c_1, c_2, c_3\}$
    \STATE \textbf{// Set target}
    \STATE $y_i \leftarrow c_t$ \quad // Label = class in target cell (top-left, $t=0$)
    \STATE \textbf{// Build grid image}
    \STATE Initialize $G_i \in \mathbb{R}^{3 \times 448 \times 448}$
    \FOR{cell $j = 0$ to $m^2-1$}
        \STATE Pop image index from pool: $\text{idx} \leftarrow P_{c_j}.\text{pop}()$
        \STATE Load image: $x \leftarrow D[\text{idx}].\text{image}$ \quad // $224 \times 224$
        \STATE Compute cell position: $(r, c) = (j \div m, j \mod m)$
        \STATE Place in grid: $G_i[:, r \cdot 224:(r+1) \cdot 224, c \cdot 224:(c+1) \cdot 224] \leftarrow x$
    \ENDFOR
    \STATE \textbf{// Create ground-truth mask}
    \STATE $M_i \in \{0,1\}^{448 \times 448}$, all zeros
    \STATE $M_i[0:224, 0:224] \leftarrow 1$ \quad // Mark target cell (top-left)
\ENDFOR
\RETURN $\{(G_i, y_i, M_i)\}_{i=1}^N$
\end{algorithmic}
\end{algorithm}

\subsection{Multi-Head Model}

Instead of training a model on grid images, we load a pre-trained single-head ResNet-18 checkpoint (trained on individual $224 \times 224$ lesion images) into a multi-head architecture, as described in Algorithm~\ref{alg:multihead_init}. This approach enables localization evaluation without requiring grid-specific training.

\begin{algorithm}[H]
\caption{Multi-Head Model Initialization with Checkpoint Reuse}
\label{alg:multihead_init}
\begin{algorithmic}[1]
\REQUIRE Pre-trained checkpoint $\theta_{\text{single}}$ (ResNet-18 on ISIC), number of heads $H=4$, number of classes $C$
\ENSURE Multi-head model $f_{\text{multi}}$ with shared backbone and $H$ classification heads
\STATE \textbf{Initialize multi-head architecture}
\STATE $f_{\text{multi}} \leftarrow$ GridMultiHead(backbone=ResNet-18, num\_heads=$H$, num\_classes=$C$)
\STATE \textbf{Load pre-trained backbone weights}
\STATE Load checkpoint: $\theta_{\text{single}} \leftarrow$ torch.load(checkpoint\_path)
\STATE Extract backbone weights: $\theta_{\text{backbone}} \leftarrow \{\text{k: v for k, v in } \theta_{\text{single}} \text{ if "backbone" in k}\}$
\STATE Update model: $f_{\text{multi}}.\text{backbone.load\_state\_dict}(\theta_{\text{backbone}})$
\STATE \textbf{Initialize all heads identically}
\STATE Extract head weights: $\theta_{\text{head}} \leftarrow \{\text{k: v for k, v in } \theta_{\text{single}} \text{ if "classifier" in k}\}$
\FOR{$h = 0$ to $H-1$}
    \STATE $f_{\text{multi}}.\text{heads}[h].\text{load\_state\_dict}(\theta_{\text{head}})$
\ENDFOR
\RETURN $f_{\text{multi}}$
\end{algorithmic}
\end{algorithm}

\subsection{DiFull Attribution Wrapper}

The DiFull wrapper (Algorithm~\ref{alg:difull_wrapper}) modifies how the model processes grid images during attribution. Unlike cropping (which was never used during our evaluation), it passes the full grid through the backbone and extracts logits from the target head. Since the full grid enters the backbone, gradients $\nabla_G z$ can flow to \emph{all} pixels. The backbone was trained on individual lesions, so its features naturally respond most strongly to lesion-like patterns. When the target head is applied to full-grid features, attribution should concentrate on the target cell (which contains a lesion) and ignore other cells. We know the correct region (top-left cell), and the GridPG metric measures whether certified pixels align with this known target.



\begin{algorithm}[H]
\caption{DiFull Attribution Wrapper}
\label{alg:difull_wrapper}
\begin{algorithmic}[1]
\REQUIRE Multi-head model $f_{\text{multi}}$, target head index $h_t=0$, full grid image $G \in \mathbb{R}^{3 \times 448 \times 448}$
\ENSURE Logits $z \in \mathbb{R}^C$ for target cell
\STATE \textbf{// Forward pass: full grid through backbone}
\STATE $\phi(G) \leftarrow f_{\text{multi}}.\text{backbone}(G)$ \quad // Global features from full $448 \times 448$ grid
\STATE \textbf{// Apply target head}
\STATE $z \leftarrow f_{\text{multi}}.\text{heads}[h_t](\phi(G))$ \quad // Logits from head 0
\RETURN $z$
\end{algorithmic}
\end{algorithm}

\subsection{Certified Grid Evaluation GridPG Metric}

The Certified GridPG score quantifies localization accuracy:
$$
\text{GridPG} = \frac{\#\{\text{certified-1 pixels in target cell}\}}{\#\{\text{all certified-1 pixels}\}}
$$
\textbf{GridPG = 1.0:} Perfect localization—all certified important pixels are in the target cell. \textbf{GridPG = 0.25:} Random baseline for $2\times2$ grid (1/4 cells). \textbf{GridPG < 0.25:} Misleading—certified pixels concentrate in \emph{wrong} cells, worse than random.

\subsection{Grid ISIC Localization Results}

Grid-based evaluation was conducted on 100 synthetic $2 \times 2$ grid images constructed from ISIC skin lesion samples. Each grid contains four different lesion classes, with the target class always placed in the top-left cell (cell 0). Certification was performed using the same randomized smoothing protocol ($\sigma=0.15$, $n=100$, $\tau=0.75$) across all attribution methods at three sparsity levels ($K=\{5\%, 25\%, 50\%\}$).

Table~\ref{tab:grid_isic_localization} presents the Certified GridPG scores for all attribution methods. Higher scores indicate better alignment between certified important pixels and the known target region.

\begin{table}[H]
\centering
\caption{Grid ISIC Localization Results: Certified GridPG Scores. Values indicate the fraction of certified important pixels that correctly localize to the target cell. Theoretical random expectation for $2\times2$ grid would be 0.25 (1/4 cells).}
\label{tab:grid_isic_localization}
\begin{tabular}{lccc}
\toprule
\textbf{Method} & \textbf{K=50\%} & \textbf{K=25\%} & \textbf{K=5\%} \\
\midrule
Occlusion            & \textbf{0.3357} & \textbf{0.3542} & 0.2937 \\
GradCAM              & 0.3067          & 0.3292          & \textbf{0.4204} \\
LRP                  & 0.2691          & 0.2654          & 0.1950 \\
IntegratedGradients  & 0.1845          & 0.0000          & 0.0000 \\
RISE                 & 0.0000          & 0.0000          & 0.0000 \\
\bottomrule
\end{tabular}
\end{table}

\paragraph{Critical Analysis of Grid Localization Performance.}
The grid-based localization results reveal significant limitations in the certified attribution framework when applied to synthetic multi-region images. Most methods achieve GridPG scores near or below the theoretical random expectation of 0.25 (for a $2\times2$ grid where each cell has equal probability), with three methods (IntegratedGradients at $K=25\%$ and $K=5\%$, RISE at all $K$ values) scoring exactly 0.0, indicating complete localization failure.

\textbf{RISE consistently fails to produce any certified important pixels in the target region.} Across all sparsity levels, RISE achieves GridPG = 0.0, meaning that none of its certified pixels align with the known ground truth location. This catastrophic failure suggests that RISE's stochastic masking procedure produces highly unstable attribution maps under randomized smoothing. The method's reliance on random occlusion patterns appears fundamentally incompatible with certification, as small input perturbations drastically alter which masks contribute to the final attribution.

\textbf{IntegratedGradients also fails at higher sparsity levels ($K=25\%$ and $K=5\%$).} While IG performed well on standard medical images in Experiments 1 and 2, it achieves 0.0 GridPG on grid images for sparse settings. This breakdown may stem from the grid structure confusing gradient-based attribution: the model's gradients distribute across multiple lesion-containing cells rather than concentrating on the target cell, and sparsification then selects pixels from incorrect regions.

\textbf{Even the best-performing methods barely exceed the theoretical random expectation.} GradCAM achieves the highest score (0.4204 at $K=5\%$), which is only moderately above the 0.25 theoretical random level (1/4 cells). Occlusion reaches 0.3542 at $K=25\%$, similarly modest. These low scores indicate that certified attributions do not reliably localize to known target regions, even when the ground truth mask is unambiguous.

\textbf{Certification visualizations exhibit extreme sparsity and abstention.} Qualitative inspection of certified attribution maps on grid images reveals that the vast majority of pixels are marked as "abstain" rather than receiving confident important/unimportant labels. The few certified pixels that do appear are scattered randomly across the grid rather than concentrating in the target cell. This suggests that the randomized smoothing procedure, when applied to grid images, produces highly inconsistent top-$K$ rankings that fail to reach the certification threshold $\tau=0.75$.

\paragraph{Potential Causes and Methodological Limitations.}
Several factors may contribute to these poor localization results:

\begin{enumerate}
\item \textbf{Model-data mismatch:} The ResNet-18 backbone was trained on individual $224 \times 224$ lesion images, not on $448 \times 448$ grids. When confronted with a grid containing four distinct lesions, the model's internal representations may not align with the assumption that only one cell is relevant. The multi-head wrapper attempts to address this by selecting the target head, but the backbone's features are still derived from the full grid, potentially diffusing attention across all four cells.

\item \textbf{DiFull wrapper limitations:} Unlike cropping-based approaches, DiFull processes the entire grid through the backbone. While this allows gradients to flow to all pixels (necessary for full-image attribution), it also means that the backbone "sees" all four lesions simultaneously. Attribution methods may then distribute importance across multiple cells, and randomized smoothing further destabilizes this already-ambiguous signal.

\item \textbf{Insufficient training on grid structure:} The checkpoint was never fine-tuned or trained on grid images. The model's weights encode priors from single-lesion classification, which may be incompatible with the grid task. This mismatch likely reduces both prediction confidence and attribution coherence on grid images.

\item \textbf{Attribution method assumptions violated:} Gradient-based methods assume a single dominant signal path from input to output. In grid images with four competing lesions, this assumption breaks down. Perturbation-based methods (RISE, Occlusion) rely on masking regions to observe prediction changes, but in grids, masking one cell may expose another lesion, confounding the importance signal.
\end{enumerate}

\textbf{Implications for faithfulness and robustness on standard medical images.} The grid evaluation casts doubt on the validity of certification results from Experiments 1 and 2. If certified attributions fail so dramatically on grid images with known ground truth, it raises the question: are the certified pixels on real medical images (ChestX-ray, Brain MRI, ISIC, APTOS) truly meaningful, or are they similarly arbitrary?

Revisiting the faithfulness curves from Experiments 1 and 2 with this skepticism reveals additional concerns. Many faithfulness curves exhibit erratic behavior—either remaining nearly flat (indicating that removing "certified important" pixels barely affects the model's confidence) or showing extreme non-monotonic oscillations (confidence increases after deleting pixels, then drops again). A proper faithfulness curve should decrease \emph{smoothly and monotonically} as more important pixels are removed. The observed irregularities suggest that:

\begin{itemize}
\item \textbf{Certified pixels may not be genuinely important.} Flat curves indicate that the model's decision does not actually depend on the certified regions, contradicting the attribution method's claims.
\item \textbf{Randomized smoothing may amplify noise rather than reveal robust structure.} The certification process, instead of filtering out unstable attributions, may be certifying random pixel patterns that happen to be consistent across noisy samples by chance.
\item \textbf{Model training quality may be insufficient.} If the underlying classifiers are not strongly confident or rely on spurious correlations, their attribution maps will be inherently unstable regardless of the certification method.
\end{itemize}

These findings suggest that pixel-level certification via randomized smoothing, while theoretically sound, may require substantially stronger model training, more careful hyperparameter tuning, or fundamentally different attribution methods to achieve practical utility in medical imaging applications.


\end{document}





